\section{闪光释放的多巴胺}
\label{sec:dishdopamine}

据我们所知，唯一一项直接让多巴胺给培养物当老师的实验，是在一个研究者可远程接入的类器官平台上完成的~\cite{Jordan2024}。在浸泡类器官的液体中加入装在光敏“笼”里的多巴胺：被笼锁住的分子不作用于受体。笼锁多巴胺的浓度为 $1$~mM。需要奖励网络时，用紫外光照射：笼分解，多巴胺释放到周围空间。

回路是这样的。同一个刺激反复施加；如果它引发的脉冲比以前多，就打开闪光、释放多巴胺。在 $240$ 次重复中，诱发脉冲数总体上在增加。作者自己写道，仅凭这一项实验不能断定增长是多巴胺引起的~\cite{Jordan2024}。

对工程师来说，这是在培养皿中搭建三因子规则~\eqref{eq:threefactor} 的第一次尝试：发送者和接收者的活动留下资格迹，由一个公共信号决定是否把变化固定下来。但这里的信号只能为正：有奖励或没有奖励，装置给不出负误差 $\delta<0$。而且多巴胺扩散和清除都很慢，就像~\eqref{eq:lowpass} 中的浓度那样，因此频繁的闪光可能只是抬高了它的总体水平。要把学习与这种效应区分开，需要两个对照：在与网络响应无关的随机时刻给同样多的闪光；以及不加笼锁多巴胺、只给同样的闪光。还需要休息后的检验：多巴胺消散后，变化是否依然存在。

\section{多巴胺教硅片}
\label{sec:biohybrid}

在反方向的实验中，活细胞教电子器件~\cite{Keene2020}。把分泌多巴胺的细胞培养在一个有机神经形态元件上方的微通道中，这个元件是一只晶体管，其沟道电导充当突触权重。细胞释放的多巴胺在元件电极处被氧化，从而改变电导。该器件还模拟了突触间隙中递质的清除机制，因此权重不仅能长期改变，还能复原。

从电路角度看，这是一个化学输入的泄漏积分器：权重累积多巴胺流量，而“清除”决定它遗忘的快慢，与~\eqref{eq:lowpass} 是同一个RC电路。这里的关键是连接方向。第~\ref{sec:debugging}~章的探针看不到广播寄存器，因为它们是化学的。这个元件读取的恰恰是化学寄存器，并把它转换成电学权重。对脑机接口来说，这是通向一种传感器的路径：它不仅能听到数据总线，还能听到控制寄存器。

\section{自带多巴胺神经元的类器官}
\label{sec:midbrainorg}

人们已经能用人类干细胞培养出类器官，它们类似于大脑中\nt{DOPA}神经元所在的那个区域。其中有能工作、能分泌多巴胺的多巴胺神经元~\cite{Jo2016}。这类类器官主要用于研究疾病。我们没有找到用它们的多巴胺给另一个网络当老师的实验。

对于用活神经元搭成的机器，这正是缺失的部件。第~\ref{sec:livingmachine}~章的试验台是没有控制寄存器的数据总线加流量控制；这里已经有了广播信号源。缺的是评论器：一个计算预测误差~\eqref{eq:ctd} 并控制这些神经元的网络。把皮层类器官和这样的类器官连起来，让多巴胺按误差释放，是一个开放问题，目前只是假说，还不是实验。

\section{不靠多巴胺平衡杆子的类器官}
\label{sec:cartpole}

近期最完整的一项实验完全不用多巴胺~\cite{Robbins2026}。把小鼠皮层类器官接入“小车倒立摆”任务的闭环：类器官的活动推动小车，杆子的位置通过刺激反馈回来。两次尝试之间，类器官接受短暂的高频教学刺激。具体给哪些刺激，由一个强化学习程序选择。

结果是测量出来的：
\begin{itemize}
\item 对大多数类器官，程序选择的教学信号比随机选择的信号或没有信号效果更好；
\item 休息 $45$ 分钟后，改进没有保留；
\item 阻断两类谷氨酸受体AMPA和NMDA，改进消失。
\end{itemize}

最后一条说明学到的东西存在哪里：在\nt{GLUT}突触中，并且经由第~\ref{sec:weightstore}~节中充当输入—输出符合检测器的那个受体。第二条说明缺了什么：有快速变化，没有巩固，就像第~\ref{sec:motorlearning}~章中只有快学习者、没有慢学习者的情形。老师的角色则变了：挑选电流模式的工程师被程序取代，但老师仍然站在培养皿之外。

在更早的电极阵列培养物学习实验中，我们也没有找到任何一项使用多巴胺的：教学信号一律是电刺激。

\section{谁在教培养物}

\begin{table}[t]
\centering
\footnotesize
\setlength{\tabcolsep}{4pt}
\renewcommand{\arraystretch}{1.15}
\begin{tabular}{>{\raggedright}p{2.7cm}>{\raggedright}p{2.5cm}>{\raggedright}p{3.2cm}>{\raggedright\arraybackslash}p{4.4cm}}
\toprule
实验 & 谁来教 & 教学信号 & 已知情况 \\
\midrule
出现所需响应时停止刺激 & 工程师 & 刺激结束 & 响应被固定——已测量 \\
DishBrain（第~\ref{sec:dishbrain}~章） & 工程师 & 可预测电流或随机电流 & 一次会话内回合数的显著增长只出现在完整反馈下，安静条件下效应较小——已测量；跨会话不稳定；原因仍有争论 \\
小车倒立摆 & 强化学习程序 & 程序选择的高频刺激 & 优于随机刺激，但休息后不保留——已测量 \\
闪光释放多巴胺 & 规则“脉冲更多即奖励” & 光释放的多巴胺 & 脉冲增加——观察结果；多巴胺的作用未证实 \\
生物混合突触 & 活细胞 & 细胞分泌的多巴胺 & 元件权重的长期改变与复原——已测量 \\
含\nt{DOPA}神经元的类器官 & --- & --- & 有多巴胺来源；未用作老师 \\
\bottomrule
\end{tabular}
\caption{谁在用什么教芯片上的活神经元。}
\label{tab:dishteachers}
\end{table}

在所有由培养物本身学习的实验中，老师目前都在外面（表~\ref{tab:dishteachers}）：工程师、程序，或工程师设定的规则。多巴胺只进过培养皿一次，它的作用还有待证明。在培养皿中搭建真正的广播通道，像第~\ref{sec:dofa}~章那样带有评论器、误差和资格迹，是下一步，而这一步还没有人迈出。
